\begin{statement}{A.28}{Proved}{The three-tail reduction and the remaining two-boundary algebra}
Let $C:v^2=w^7-w$, with its two cusps $P_0,P_\infty$.
After a finite extension at the place above two, its stable pointed
special fibre has three elliptic components attached to one rational
component carrying both cusps. The three elliptic components map
isomorphically to the good special fibre of $E_{0,-1}$ under the
cyclic cubic covering
\[
                C\longrightarrow E_{0,-1},\qquad
                      (w,v)\longmapsto(w^3,wv).
                                                        \tag{A.28.1}
\]
The cubic rotation permutes the three tails, whose edge lengths are
equal and strictly positive.

For every odd $\ell$ and every original packet, the monodromy of a
local Frobenius--tame commutator has the following alternatives on
the free rational graded Lie algebra
\[
           \mathbb L(kt\oplus V_1\oplus V_2\oplus V_3),
           \qquad \dim V_i=2,\quad k=\mathbb Q_\ell.
\]
Here $t$ is the primitive class of the selected cusp. If its first
Johnson derivation is nonzero, the actual fixed subgroup is precisely
that cusp inertia. Otherwise its second Johnson derivation, together
with one actual CM Frobenius conjugate, has common kernel
\[
                         \mathbb L_k(t,\Omega),
              \qquad \Omega=\omega_1+\omega_2+\omega_3,
                                                        \tag{A.28.2}
\]
where $\omega_i$ is the symplectic form of $V_i$.
Consequently every leading graded class of an actual fixed element
lies in this displayed algebra. The algebra contains higher
commutators, so this statement leaves the corresponding exceptional
full fixed subgroup undecided.
\end{statement}
\begin{proof}
We first determine the geometry, keeping the distinction between the
cyclic cubic quotient and the degree-six map to the original $E$.
On $E'=E_{0,-1}$, the function $x$ has divisor $2T-2O$, with
$T=(0,0)$ of order two. The cubic cover is the Kummer cover $w^3=x$.
Both branch points specialize to the origin after potential good
reduction above two: the special fibre is supersingular, and its
geometric two-torsion consists only of the origin. The stable
two-pointed model of this good elliptic curve therefore consists of
the elliptic component and a rational component bearing $T$ and $O$,
joined by an edge of positive length. One obtains this model by
successively blowing up the coincidence of the two sections, then
contracting unmarked valence-two rational components.

Normalize this marked model in $w^3=x$, after a finite extension
which makes the vertical Kummer valuations divisible by three.
This is a tame operation, since the residue characteristic is two.
The local description by a cube root of a monomial in the parameters
shows that its only possible branch points on each component are the
marked points and the node. At the elliptic component the node is the
only such point. A cyclic cubic cover cannot branch at just one
point: the divisor of a Kummer function has degree zero, so its sole
possibly nonzero order modulo three must be zero. Thus the cover is
unramified at this node.

Over the rational component the two marked orders are $2$ and $-2$
modulo three. Their sum is zero, again giving no ramification at
the node. Its normalization is a connected rational curve, by
Riemann--Hurwitz for a cyclic cubic cover of a projective line with
exactly two totally ramified points. There are three points over its
node. Over the elliptic component there are exactly two possibilities:
a connected etale degree-three isogeny, or three disjoint copies of
that elliptic curve. These statements also verify semistability
directly: the ramified marked-point coordinates have the form
$u=z^3$, and each unramified node has three ordinary nodal preimages.
After the preceding finite extension there are no further positive-
genus components introduced by normalization or regular subdivision.

The first possibility would give one elliptic component and one
rational component joined at three nodes, hence a dual graph of
first Betti number two. It is impossible. Indeed A.4 gives
$J_C\sim E^3$ after a finite extension; its good reduction then
follows from that of $E$ and Neron--Ogg--Shafarevich. The torus rank
of the generalized Jacobian of a nodal fibre is the first Betti
number of its dual graph, so that graph must be a tree. We obtain
the asserted three-tail fibre. The rotation fixes the central
component, permutes its three nodes, and permutes the elliptic
components. It preserves the node thicknesses, proving equality
and positivity of the three edge lengths. Both cusps lie on the
central rational component, at the totally ramified marked points.

Choose an ordinary base point on that central component, after a
finite extension. The affine edge-twist theorem used in A.27 gives
zero first Johnson degree and second-degree derivation
\[
 D_0(t)=0,\qquad D_0|_{V_i}=m_i\operatorname{ad}(\omega_i),
                     \qquad m_i\ne0.                   \tag{A.28.3}
\]
For a Frobenius--tame commutator, its coefficients are the positive
edge lengths multiplied by the common nonzero scalar $q-1$.
The argument needs only their nonvanishing.

Compare this affine representative with the actual packet action.
Their difference is inner. Both actions fix the leading selected
cusp class $t$. The first-degree inner term is therefore
$\lambda\operatorname{ad}(t)$: the centralizer of a free generator
in the free Lie algebra is its own line. If $\lambda\ne0$, its
kernel is $kt$. Primitivity of the actual cusp class, subtraction
of its actual power, and separation of the lower-central filtration
give the full fixed subgroup assertion, just as in A.27.
If $\lambda=0$, the first-degree action vanishes. The second-degree
inner correction is $\operatorname{ad}(z_2)$ and must kill $t$,
so $z_2=0$. Thus (A.28.3) is the actual second-degree derivation,
with no unspecified inner term.

The first-core normal form of A.27 applies to any number of elliptic
blocks and an additional free generator $t$. Treat $t$ and the
$\omega_i$ as the boundary-prefix letters. On a first nonboundary
core of type $i$, the only term increasing its prefix is nonzero
right multiplication by $m_i\omega_i$. Hence the associative
kernel is exactly $k\langle t,\omega_1,\omega_2,\omega_3\rangle$;
its primitive part is the corresponding free Lie algebra.

We compute its intersection with an actual arithmetic conjugate.
Extend coefficients to split the CM field, and choose symplectic
bases $a_i,b_i$ on the tails with $(f_i)_*a_i=x$ and
$(f_i)_*b_i=y$ on $E'$, where $x,y$ is a CM eigenbasis. This is
possible because each component map has degree one. Put
$A=\sum a_i$ and $B=\sum b_i$.
The original elliptic subspace is $\langle A,B\rangle$ and its
orthogonal complement is the sum-zero subspace in the two triples.
The automorphism $(w,v)\mapsto(-w,iv)$ acts as $[-i]$ on the
original elliptic factor and as $[i]$ on the complementary two
factors: on the three holomorphic differentials its eigenvalues
are $i,-i,i$. It preserves each special elliptic tail, since it
has order four and commutes with their three-cycle. On every tail
its action is identified by (A.28.1) with $[-i]$ on $E'$.
It follows that the global CM eigenspaces are
\[
          \langle A\rangle\oplus\{\textstyle\sum u_i b_i:
                        \sum u_i=0\},\qquad
          \langle B\rangle\oplus\{\textstyle\sum u_i a_i:
                        \sum u_i=0\}.                  \tag{A.28.4}
\]

An actual good ordinary Frobenius with the two CM eigenvalues
$s,r$ acts, in these coordinates, by
\[
 a_i\longmapsto r a_i+(s-r)A/3,\qquad
 b_i\longmapsto s b_i+(r-s)B/3.                         \tag{A.28.5}
\]
Take its place away from $\ell$ and pass to any required positive
residue power. Its eigenvalue ratio is not a root of unity and has
complex absolute value one. Hence $s-r\ne0$ and $s+2r\ne0$.
Put $P=\frac13(1)_{3\times3}$ and
$M=rI+(s-r)P$. The other matrix is $srM^{-1}$.
The plane $U=\langle\omega_1,\omega_2,\omega_3\rangle$
identifies with the diagonal matrices in the $a_i\wedge b_j$
coordinates. The conjugate plane is $srMUM^{-1}$.
Every entry of $M$ is nonzero. If a diagonal matrix $D$ satisfies
$DM=MD'$ for another diagonal matrix $D'$, the $(i,j)$ entry gives
$D_{ii}=D'_{jj}$ for every $i,j$. Both are therefore scalar. Thus
\[
                             U\cap gU=k\Omega.         \tag{A.28.6}
\]
The exact symbolic matrix calculation is also executed in
\texttt{a\_C3\_stable\_graph.py}.

The boundary generator $t$ is carried to a nonzero scalar multiple
of itself. Decompose tensor words according to the occurrences of
$t$ and then group each intervening even-length $V$ block into
pairs. The tensor-intersection argument of A.27 on each block
shows that the two associative kernel algebras intersect in
$k\langle t,\Omega\rangle$. Their primitive intersection is
(A.28.2). The actual commutator group is normal in a sufficiently
small arithmetic group, so the Frobenius conjugate used here is
another actual derivation fixing every element under consideration.
Finally descend the kernel statement to $\mathbb Q_\ell$.

Unlike the one-cusp case, the free Lie algebra in (A.28.2) is not
one-dimensional. In particular $[t,\Omega]$ survives this argument.
The three genus-two quotient maps do not remove it at leading
weight, because each sends the two generators to proportional
symplectic forms. A further actual cusp-path or higher-monodromy
calculation is required; no existence of a fixed element with
this leading term is being asserted.
\end{proof}
